\begin{examplebox}
	\textbf{\textcolor{CExample}{例 1（基础）}}
	\textbf{\textcolor{CExample}{题目：}}
	在 $\triangle ABC$ 中，已知 $A(0,0)$，$B(4,0)$，$C(1,3)$，求垂心 $H$ 的坐标。
	\textbf{\textcolor{CExample}{解题步骤：}}
	\begin{description}[style=nextline, leftmargin=2.8em, labelsep=0.5em, itemsep=0.45em]
		\item[\textbf{第一步（设未知量）：}] 设 $H(x,y)$。
			\par\textbf{理由：} 垂心坐标未知，待求。
		\item[\textbf{第二步（列向量方程）：}] 由垂心垂直性质，$AH\perp BC$，$BH\perp AC$。计算：
			\[
				\overrightarrow{AH}=(x,y),\quad \overrightarrow{BC}=(-3,3),\quad \overrightarrow{BH}=(x-4,y),\quad \overrightarrow{AC}=(1,3).
			\]
			由 $\overrightarrow{AH}\cdot\overrightarrow{BC}=0$ 得 $-3x+3y=0$，即 $x-y=0$；由 $\overrightarrow{BH}\cdot\overrightarrow{AC}=0$ 得 $x+3y-4=0$。
			\par\textbf{理由：} 垂直 $\Rightarrow$ 向量内积为0。
		\item[\textbf{第三步（解方程组）：}] 解方程组
			\[
				\begin{cases}
					x-y=0,\\
					x+3y-4=0,
			\end{cases}
			\]
			得 $x=1$，$y=1$。所以 $H(1,1)$。
			\par\textbf{理由：} 二元一次方程组求解。
	\end{description}
	\textbf{\textcolor{CExample}{关键结论：}}
	\textbf{垂心 $H$ 的坐标为 $(1,1)$。}
	\medskip
	\textbf{\textcolor{CExample}{例 2（稍变形）}}
	\textbf{\textcolor{CExample}{题目：}}
	在锐角 $\triangle ABC$ 中，$H$ 是垂心，若 $\angle A=60^\circ$，求 $\angle BHC$ 的大小。
	\textbf{\textcolor{CExample}{解题步骤：}}
	\begin{description}[style=nextline, leftmargin=2.8em, labelsep=0.5em, itemsep=0.45em]
		\item[\textbf{第一步（找垂直关系）：}] 过 $B$ 作 $AC$ 垂线，垂足为 $E$；过 $C$ 作 $AB$ 垂线，垂足为 $F$。则 $H$ 为 $BE$ 与 $CF$ 的交点，故 $\angle AEH=\angle AFH=90^\circ$。
			\par\textbf{理由：} 垂心定义及垂直性质。
		\item[\textbf{第二步（四边形内角和）：}] 四边形 $AEHF$ 中，$\angle AEH+\angle AFH=180^\circ$，故 $\angle A+\angle EHF=180^\circ$。
			\par\textbf{理由：} 四边形内角和 $360^\circ$。
		\item[\textbf{第三步（对顶角转化）：}] $\angle BHC=\angle EHF$（对顶角），所以 $\angle BHC=180^\circ-\angle A=120^\circ$。
			\par\textbf{理由：} 对顶角相等。
	\end{description}
	\textbf{\textcolor{CExample}{关键结论：}}
	\textbf{$\angle BHC=120^\circ$。}
	\medskip
	\textbf{\textcolor{CExample}{例 3（综合应用）}}
	\textbf{\textcolor{CExample}{题目：}}
	已知 $H$ 是 $\triangle ABC$ 的垂心，求证：
	\[
		\overrightarrow{AH}\cdot\overrightarrow{BC}+\overrightarrow{BH}\cdot\overrightarrow{AC}+\overrightarrow{CH}\cdot\overrightarrow{AB}=0.
	\]
	\textbf{\textcolor{CExample}{解题步骤：}}
	\begin{description}[style=nextline, leftmargin=2.8em, labelsep=0.5em, itemsep=0.45em]
		\item[\textbf{第一步（单个垂直）：}] 由 $H$ 为垂心，得 $\overrightarrow{AH}\cdot\overrightarrow{BC}=0$。
			\par\textbf{理由：} $AH\perp BC$，点积为零。
		\item[\textbf{第二步（另外两个垂直）：}] 同理，$\overrightarrow{BH}\cdot\overrightarrow{AC}=0$，$\overrightarrow{CH}\cdot\overrightarrow{AB}=0$。
			\par\textbf{理由：} $BH\perp AC$，$CH\perp AB$，点积为零。
		\item[\textbf{第三步（求和得证）：}] 三式相加得
			\[
				\overrightarrow{AH}\cdot\overrightarrow{BC}+\overrightarrow{BH}\cdot\overrightarrow{AC}+\overrightarrow{CH}\cdot\overrightarrow{AB}=0.
			\]
			\par\textbf{理由：} 等式相加。
	\end{description}
	\textbf{\textcolor{CExample}{关键结论：}}
	\[
		\overrightarrow{AH}\cdot\overrightarrow{BC}+\overrightarrow{BH}\cdot\overrightarrow{AC}+\overrightarrow{CH}\cdot\overrightarrow{AB}=0 \text{ 成立.}
	\]
\end{examplebox}
